\subsubsection{Vormen van uitlegbaarheid}
\label{subsubsec:vormen-van-uitlegbaarheid}

De literatuur kent geen eenduidige definitie voor \emph{interpretability} \autocite{Lipton2016}. In deze scriptie wordt \emph{interpretability}, of uitlegbaarheid, opgevat als de mate waarin een systeem de redenen voor zijn beoordeling kan weergeven op een manier die een mens kan volgen. Op basis van inzichten uit de sociale wetenschappen stelt \textcite{Miller2017} dat mensen daarbij geen volledige opsomming van oorzaken verwachten, maar enkel de belangrijkste factoren die tot een resultaat bijdragen. Voor dit onderzoek luidt de vraag die een uitleg moet beantwoorden dus: waarom beoordeelt het systeem deze e-mail als risicovol?

Hoe een uitleg tot stand komt, hangt rechtstreeks samen met het type model. Neem bijvoorbeeld een interpreteerbaar model dat op regels gebaseerd is, zoals besproken in sectie~\ref{subsubsec:rulebased-detectie}. De interne werking van zo'n model is goed te volgen en te reproduceren: elke regel is een expliciete test met een binaire uitkomst, en aan elke regel kan vooraf een uitleg worden gekoppeld die getoond wordt wanneer de regel waar is. Bij de \emph{machine learning}-methoden uit sectie~\ref{subsubsec:machinelearning-detectie} is dat veel moeilijker af te dwingen, waardoor de uitleg achteraf moet worden opgebouwd. De literatuur noemt dit \emph{post-hoc} uitleg. Twee veelgebruikte technieken zijn LIME \autocite{Ribeiro2016}, dat het gedrag van het model rond één voorspelling benadert met een eenvoudiger model, en SHAP \autocite{Lundberg2017}, dat berekent in welke mate elk kenmerk bijdraagt aan een voorspelling. Post-hoc uitleg is echter niet zonder risico: beide technieken benaderen het model en blijken zich te laten misleiden, waardoor de uitleg een gunstiger beeld kan schetsen dan het model verdient \autocite{Slack2019}. \textcite{Rudin2019} pleit er daarom voor om interpreteerbare modellen te verkiezen, zeker naarmate de impact van een beslissing groter wordt.

Verder kan er een onderscheid worden gemaakt tussen globale en lokale uitleg. Globale uitleg geeft inzicht in de werking van een model als geheel, bijvoorbeeld door aan te tonen welke kenmerken een belangrijke rol spelen bij het detecteren van phishingmails. Lokale uitleg richt zich daarentegen op een individuele voorspelling en verduidelijkt welke kenmerken hebben bijgedragen aan de risicoscore van een specifieke e-mail \autocite{Guidotti2018}.